% GENERATED FILE. Edit canonical lessons, then run npm run generate:books.
% canonical-source-hash: fnv1a64:75ae71548ea388c8
% canonical-lessons: ZH-C258-youhao, ZH-C258-rexin, ZH-W258-re-hot, ZH-C258-naixin, ZH-C258-youmo, ZH-C258-xingge

\chapter{Warm-Hearted}
\label{ch:zh-warm-hearted}
\hlchaptermodality{258}

\begin{chapteropening}
\textbf{By the end of this chapter:} \emph{I can say that someone is friendly, warm-hearted, patient or witty, and talk about their personality.}

Complete the last lesson of chapter 258: I can say that someone is friendly, warm-hearted, patient or witty, and talk about their personality.
\end{chapteropening}

\section[yǒuhǎo]{yǒuhǎo — friendly}

\label{lesson:ZH-C258-youhao}



\begin{quote}
Before the new one: say the Mandarin for a singer, then the Mandarin for to applaud.
\end{quote}

\subsection*{You'll want to know: yǒuhǎo}
\emph{yǒuhǎo} — \textquotedblleft{}friendly\textquotedblright{}.

\begin{quote}
\emph{zhèr de rén hěn yǒuhǎo} — People here are very friendly.
\end{quote}

Its characters come later; for now it is a word for the ear.

Two third tones meet in this word: yǒu is said with a rising tone before hǎo, and the spelling keeps both marks.

\begin{cousinweb}
\emph{yǒuhǎo} is built from \emph{yǒu} 'a friend' and \zh{好} 'good', both inherited from Middle Chinese, \emph{hjuwX} and \emph{xawX} (Guangyun, in Baxter's transcription), which glosses the first 'a friend'.
\end{cousinweb}

\begin{grammarlens}[title={What you've built}]
One more word for describing things.
\end{grammarlens}

\subsection*{Your turn}
\begin{itemize}
\raggedright
  \item \emph{Say it:} \emph{yǒuhǎo}
  \item \emph{Say it:} \emph{yǒuhǎo}, once more
  \item \emph{Say it:} \emph{zhèr de rén hěn yǒuhǎo} — People here are very friendly.
  \item \emph{Recall:} say the Mandarin for a singer, then the Mandarin for to applaud, then say \emph{yǒuhǎo} again
\end{itemize}

\begin{tcolorbox}[breakable,colback=teal!4,colframe=teal!35!black,title={Before you move on}]
What does yǒuhǎo mean? (\textquotedblleft{}Friendly\textquotedblright{}.) Say it once more.
\end{tcolorbox}
\section[\texorpdfstring{\zh{热心} (rèxīn)}{热心 (rèxīn)}]{\zh{热心} (rèxīn) — warm-hearted; eager to help}

\label{lesson:ZH-C258-rexin}



\begin{quote}
Before the new one: say the Mandarin for to applaud, then the Mandarin for friendly.
\end{quote}

\subsection*{You'll want to know: \zh{热心}}
\textbf{\zh{热心}} — \emph{rèxīn} — \textquotedblleft{}warm-hearted; eager to help\textquotedblright{}.

\begin{quote}
\textbf{\zh{我们的校长很热心}} \emph{wǒmen de xiàozhǎng hěn rèxīn} — Our head teacher is very warm-hearted.
\end{quote}

The new character, in pieces: \zh{热} — ten strokes: the hand piece of \zh{打}, three strokes; beside it a stroke falling to the left, a bar that turns down and hooks up, and a dot; then four dots along the foot. The next lesson puts it in your hand.

\begin{cousinweb}
\textbf{\zh{热心}} (\emph{rèxīn}) is built from \zh{热} 'hot' and \zh{心} 'the heart', both inherited from Middle Chinese, \emph{nyet} and \emph{sim} (Guangyun, in Baxter's transcription), which glosses the first as burning, as with fire.
\end{cousinweb}

\begin{grammarlens}[title={What you've built}]
One more word for describing things.
\end{grammarlens}

\subsection*{Your turn}
\begin{itemize}
\raggedright
  \item \emph{Say it:} \emph{rèxīn}
  \item \emph{Say it:} \emph{rèxīn}, once more
  \item \emph{Say it:} \emph{wǒmen de xiàozhǎng hěn rèxīn} — Our head teacher is very warm-hearted.
  \item \emph{Recall:} say the Mandarin for to applaud, then the Mandarin for friendly, then say \emph{rèxīn} again
\end{itemize}

\begin{tcolorbox}[breakable,colback=teal!4,colframe=teal!35!black,title={Before you move on}]
What does \zh{热心} mean? (\textquotedblleft{}Warm-hearted; eager to help\textquotedblright{}.) Say it once more.
\end{tcolorbox}
\section[\texorpdfstring{\zh{热} (rè)}{热 (rè)}]{\zh{热} — the hand piece and a hooked bend, over four dots}

\label{lesson:ZH-W258-re-hot}



\begin{quote}
The word first, then the character inside it.

\begin{itemize}
\raggedright
  \item \emph{Recall:} say \emph{rèxīn} — \textbf{R1}, one lesson back, before it decays
\end{itemize}

\begin{quote}
\zh{打} \zh{心}
\end{quote}

\emph{Read it:} each one aloud, with nothing said first — \emph{dǎ}, \emph{xīn}
\end{quote}

\subsection*{Script — the hand piece and a hooked bend, over four dots}
\hlblockfigure{\detokenize{figures/ZH-W258-re-hot-filmstrip.pdf}}{How \zh{热} is written, stroke by stroke}

\begin{quote}
\zh{热}
\end{quote}

\textbf{Three pieces: two side by side on top, and a row of dots under them.}

\textbf{The top left piece, three strokes,} the hand piece of \zh{打}: a short bar, a vertical with a hook, then a stroke rising to the right across it.

\textbf{The top right piece, three strokes.} A stroke falling to the left. A bar that runs right, turns down, bends out to the right and hooks up. Then a dot inside.

\textbf{The foot, four dots,} left to right: the first leans to the left, the other three to the right.

3 + 3 + 4 = ten.

\subsection*{Your turn}
\begin{itemize}
\raggedright
  \item \emph{Write it:} \zh{热}, all ten, with the model in front of you
  \item \emph{Cover:} write it again from memory, and count the pieces, not the strokes
\end{itemize}

\begin{tcolorbox}[breakable,colback=teal!4,colframe=teal!35!black,title={Before you move on}]
How many strokes in \zh{热}? (\textbf{Ten}.) Which piece comes first? (\textbf{The hand piece of \zh{打}}.) What runs along the foot? (\textbf{Four dots}.)

Source: \href{https://raw.githubusercontent.com/chanind/hanzi-writer-data/68d10a4b21150cae5e1ebbd223eed289cf32d90c/data/\%E7\%83\%AD.json}{Hanzi Writer Data \zh{热}, pinned stroke paths and medians}.
\end{tcolorbox}
\section[nàixīn]{nàixīn — patient with people}

\label{lesson:ZH-C258-naixin}



\begin{quote}
Before the new one: say the Mandarin for friendly, then the Mandarin for eager to help.
\end{quote}

\subsection*{You'll want to know: nàixīn}
\emph{nàixīn} — \textquotedblleft{}patient with people\textquotedblright{}.

\begin{quote}
\emph{lǎoshī hěn nàixīn} — The teacher is very patient.
\end{quote}

Its characters come later; for now it is a word for the ear.

\begin{cousinweb}
\emph{nàixīn} is built from \emph{nài} 'to endure' and \zh{心} 'the heart', both inherited from Middle Chinese, \emph{nojH} and \emph{sim} (Guangyun, in Baxter's transcription), which glosses the first 'to bear'.
\end{cousinweb}

\begin{grammarlens}[title={What you've built}]
One more word for describing things.
\end{grammarlens}

\subsection*{Your turn}
\begin{itemize}
\raggedright
  \item \emph{Say it:} \emph{nàixīn}
  \item \emph{Say it:} \emph{nàixīn}, once more
  \item \emph{Say it:} \emph{lǎoshī hěn nàixīn} — The teacher is very patient.
  \item \emph{Recall:} say the Mandarin for friendly, then the Mandarin for eager to help, then say \emph{nàixīn} again
\end{itemize}

\begin{tcolorbox}[breakable,colback=teal!4,colframe=teal!35!black,title={Before you move on}]
What does nàixīn mean? (\textquotedblleft{}Patient with people\textquotedblright{}.) Say it once more.
\end{tcolorbox}
\section[yōumò]{yōumò — humorous; witty (of a person)}

\label{lesson:ZH-C258-youmo}



\begin{quote}
Before the new one: say the Mandarin for eager to help, then the Mandarin for patient with people.
\end{quote}

\subsection*{You'll want to know: yōumò}
\emph{yōumò} — \textquotedblleft{}humorous; witty (of a person)\textquotedblright{}.

\begin{quote}
\emph{wǒ bàba hěn yōumò} — My dad has a great sense of humour.
\end{quote}

Its characters come later; for now it is a word for the ear.

\begin{cousinweb}
\emph{yōumò} is usually traced to English \emph{humour}, written with two characters whose sounds come close to it, \emph{yōu} 'deep' and \emph{mò} 'silent'; who first wrote it this way is not documented in what we can open.
\end{cousinweb}

\begin{grammarlens}[title={What you've built}]
One more word for describing things.
\end{grammarlens}

\subsection*{Your turn}
\begin{itemize}
\raggedright
  \item \emph{Say it:} \emph{yōumò}
  \item \emph{Say it:} \emph{yōumò}, once more
  \item \emph{Say it:} \emph{wǒ bàba hěn yōumò} — My dad has a great sense of humour.
  \item \emph{Recall:} say the Mandarin for eager to help, then the Mandarin for patient with people, then say \emph{yōumò} again
\end{itemize}

\begin{tcolorbox}[breakable,colback=teal!4,colframe=teal!35!black,title={Before you move on}]
What does yōumò mean? (\textquotedblleft{}Humorous; witty (of a person)\textquotedblright{}.) Say it once more.
\end{tcolorbox}
\section[xìnggé]{xìnggé — personality; character}

\label{lesson:ZH-C258-xingge}



\begin{quote}
Before the new one: say the Mandarin for patient with people, then the Mandarin for witty.
\end{quote}

\subsection*{You'll want to know: xìnggé}
\emph{xìnggé} — \textquotedblleft{}personality; character\textquotedblright{}.

\begin{quote}
\emph{wǒ xǐhuan tā de xìnggé} — I like his personality.
\end{quote}

Its characters come later; for now it is a word for the ear.

\begin{cousinweb}
\emph{xìnggé} is built from \emph{xìng} 'nature' and \emph{gé} 'a pattern', both inherited from Middle Chinese, \emph{sjengH} and \emph{kaek} (Guangyun, in Baxter's transcription), which glosses the first 'nature and conduct' and the second 'a model' and 'a measure'.
\end{cousinweb}

\begin{grammarlens}[title={What you've built}]
One more word for describing things.
\end{grammarlens}

\subsection*{Your turn}
\begin{itemize}
\raggedright
  \item \emph{Say it:} \emph{xìnggé}
  \item \emph{Say it:} \emph{xìnggé}, once more
  \item \emph{Say it:} \emph{wǒ xǐhuan tā de xìnggé} — I like his personality.
  \item \emph{Recall:} say the Mandarin for patient with people, then the Mandarin for witty, then say \emph{xìnggé} again
\end{itemize}

\begin{tcolorbox}[breakable,colback=teal!4,colframe=teal!35!black,title={Before you move on}]
What does xìnggé mean? (\textquotedblleft{}Personality; character\textquotedblright{}.) Say it once more.
\end{tcolorbox}
